\documentclass[../../../include/open-logic-section]{subfiles}

\begin{document}

\olfileid{sth}{z}{story}
\olsection{Gambaran kang Luwih Rinci}
	
Ing \olref[story][approach]{sec}, kita ngutip panjelasane
Shoenfield babagan proses pambentukan himpunan. Saiki kita
arep nulis sawetara prinsip maneh supaya gambaran iki
luwih trep. Iki prinsip-prinsipe:
\begin{enumerate}
\item[] \stageshier. Saben himpunan kawangun ing sawenehing tataran.
\item[] \stagesord. Tataran-tataran diurutake: sawetara tataran
ana \emph{sadurunge} liyane.\footnote{Nyatane kita bakal nganggep
kanthi meneng-meneng yen tataran-tataran iku \emph{diurutake kanthi
becik}. Tegese diterangake ing \olref[ordinals][]{chap}.
Iki asumsi kang wigati. Nyatane, kanthi teknik kang pinter
saka \citet{Scott1974}, asumsi iki bisa \emph{diendhani}
banjur \emph{diturunake}. (Iki uga bakal nerangake yagene
kita nganggep ana tataran wiwitan.) Kita ora bisa ngrembug
iku ing kene; kanggo katrangan luwih, delengen
\citet{ButtonLT1}.}
\item[] \stagesacc. Kanggo saben tataran $S$ lan saben himpunan
kang kawangun \emph{sadurunge} tataran~$S$: ing tataran~$S$
kawangun sawijining himpunan kang anggotane persis
himpunan-himpunan kasebut. Ora ana liyane kang kawangun
ing tataran~$S$.
\end{enumerate}
Iki prinsip-prinsip informal, nanging kita bisa nggunakake
kanggo ndhukung sawetara aksioma teori himpunan Zermelo.

(Kita kudu menehi pepeling. Sanajan mengko kita ngaturake
aksioma-aksioma kang baku kanthi jeneng kang baku, prinsip-prinsip
mawa aksara miring kang anyar wae kita aturake ora nduweni
jeneng khusus ing pustaka. Kita mung menehi sebutan kang
muga-muga migunani.)

\end{document}
